% ---------------------------------------------------------------------------
%  Frontespizio
% ---------------------------------------------------------------------------
\begin{titlepage}
\thispagestyle{empty}
\begin{tikzpicture}[remember picture,overlay]
  \fill[prgsteel] (current page.north west) rectangle ([yshift=-5.5cm]current page.north east);
  \fill[prgaccent] ([yshift=-5.5cm]current page.north west) rectangle ([yshift=-5.85cm]current page.north east);
  % the worm of Verm1 on a small retro screen
  \begin{scope}
    \fill[prgink] ([xshift=1.4cm,yshift=-1.1cm]current page.north west) rectangle ([xshift=4.0cm,yshift=-3.0cm]current page.north west);
    \foreach \i in {0,...,5}{
      \fill[prggreen] ([xshift=1.60cm+\i*0.26cm,yshift=-2.60cm]current page.north west) rectangle ++(0.22cm,0.22cm);
    }
    \foreach \j in {1,...,2}{
      \fill[prggreen] ([xshift=2.90cm,yshift=-2.60cm+\j*0.26cm]current page.north west) rectangle ++(0.22cm,0.22cm);
    }
    \fill[green!80!black] ([xshift=2.90cm,yshift=-1.82cm]current page.north west) rectangle ++(0.22cm,0.22cm);
    \fill[prgred] ([xshift=3.42cm,yshift=-1.56cm]current page.north west) rectangle ++(0.22cm,0.22cm);
  \end{scope}
  \node[anchor=north west,text=white,font=\sffamily\large] at ([xshift=4.6cm,yshift=-1.5cm]current page.north west){ARC12 \textbullet\ Architettura dei Calcolatori e Laboratorio};
  \node[anchor=north west,text=white,font=\sffamily\Large\bfseries] at ([xshift=4.6cm,yshift=-2.2cm]current page.north west){RISC-V sulla piattaforma PRG32};
  \node[anchor=north west,text=prgpaper,font=\sffamily\small] at ([xshift=4.6cm,yshift=-3.4cm]current page.north west){Università degli Studi di Napoli «Parthenope» \quad$\bullet$\quad A.A. 2026/27};
\end{tikzpicture}

\vspace*{6.5cm}

\begin{flushleft}
{\sffamily\bfseries\fontsize{30}{34}\selectfont\color{prgink}
Architettura dei calcolatori}\\[6pt]
{\sffamily\fontsize{20}{24}\selectfont\color{prgsteel}
Prima programmare, poi formalizzare}\\[18pt]
{\sffamily\large\color{prgaccent}
Dalla prima istruzione \riscv a un gioco pubblicato:\\
informazione, logica, memoria e processore, insegnati sulla piattaforma PRG32\\
attraverso \vermone, il progetto che accompagna tutto il corso}
\end{flushleft}

\vfill

\begin{flushleft}
{\sffamily\Large\color{prgink} Raffaele Montella}\\[4pt]
{\sffamily\normalsize\color{prgink}
Dipartimento di Scienze e Tecnologie\\
Università degli Studi di Napoli «Parthenope»}\\[14pt]
{\sffamily\normalsize\color{prgsteel}
Implementazione di riferimento di \vermone: contributo di Giuseppe Coviello}\\[14pt]
{\small Libro di testo delle lezioni di \emph{Architettura dei Calcolatori e Laboratorio
di Architettura dei Calcolatori} (ARC12),\\ Corso di Laurea in Informatica,
primo anno, primo semestre.\\[6pt]
I listati assembly di questo libro sono letti da file che vengono assemblati\\ e provati automaticamente.}
\end{flushleft}

\vspace{1.2cm}
{\footnotesize\color{prgcomment} Piattaforma: \url{https://github.com/riscv-prg32/PRG32}\\
Host desktop e debugger: \url{https://github.com/riscv-prg32/PRG32-QT}}
\end{titlepage}
